\bolumsuz{KAYNAKLAR}
Kaynaklar metin içinde yazar soyadı ve tarih belirtilerek verilir. Kaynaklar sayfasında yazar soyadına göre alfabetik olarak sıralama yapılır. Metin içinde kaynak, cümlenin başlangıcında veya içinde verilecekse, Baran (2013) şeklinde, kaynak cümle sonunda verilecekse (Baran, 2013). şeklinde gösterilir. Nokta işareti kaynaktan hemen sonra konulur.

Kaynak birden fazla yazara ait olduğunda, yazar sayısıi ise, cümle başında veya içinde Uyar ve Yılmaz (2020) şeklinde, cümle sonunda ise (Uyar ve Yılmaz, 2020). şeklinde yazılır. Yazar sayısı ikiden fazla ise cümle başında veya içinde Uyar ve ark. (2018) şeklinde, cümle sonunda ise (Uyar ve ark, 2018). şeklinde yazılır.

Aynı yazara ait ve aynı yıl içinde yayımlanmış yayınlar Uyar (2008a), Uyar (2008b) şeklinde numaralandırılır. Aynı parantez içerisinde aynı yazarın 2 ve daha fazla eserine atıfta bulunma; son yayınlanan eseri en son belirterek aynı parantez içerisinde gösterilebilirler. Örneğin; Son araştırmalar \ldots{}. göstermiştir (Uyar, 2008, 2020).

Aynı parantez içerisinde 2 ya da daha fazla esere atıf yazar soyadlarına göre alfabetik sıraya uygun şekilde verilir; (Ateş, 2021; Çelik, 2012; Yılmaz, 2022).

İkincil kaynak durumunda (atıf yapılan kaynak başka bir kaynağa atıfta bulunuyorsa), metinde orijinal kaynağa atıfta bulunulur ve parantez içerisinde orijinal kaynağa atıfta bulunan yazara gönderme yapılır. Referans listesinde sadece orijinal kaynağa atıfta bulunan kaynak için giriş yapılır; orijinal kaynak için referans girişi yapılmaz. Kaynaklar bölümünde künyeler yazar soyadına göre sıralanır.

Başka bir kaynaktan alıntı yapma (İkincil Kaynaklar; Aktaran):

İkincil kaynak, ilgili kaynağa ulaşılamadığı durumlarda “aktaran” ifadesi ile yapılmalıdır. Örneğin;

Parantez içerisinde kullanımı:

(Baran, 1945; akt. Çoruhlu, 2015).

Metin içi kullanımı:

Baran (aktaran Çoruhlu, 2015) 1945’ te \ldots{}

\textbf{APA 7  formatına göre kaynakça yazım örnekleri}

\kaynak{\notrenk{0070C0}{\textbf{KİTAP}}}

{\teksatir \textbf{Kitap: Tek yazarlı }\par}

\kaynak{\notrenk{FF0000}{Yazarın soyadı, Yazarın adının baş harfi. (Yıl). }\notrenk{FF0000}{\textit{Kitabın adı }}\notrenk{FF0000}{(Baskı Sayısı). Yayınevi.}}

\kaynak{[Metin içinde atıf: (Jackson, 2019; Sapolsky, 2017; Tarhan, 2012]}

\kaynak{Jackson, L. M. (2019). \textit{The psychology of prejudice: From attitudes to social action} (2nd ed.). American Psychological. Association. https://doi.org/10.1037/0000168-000}

\kaynak{Sapolsky, R. M. (2017). \textit{Behave: The biology of humans at our best and worst}. Penguin Books.}

\kaynak{Tarhan, N. (2012). \textit{Mesnevi terapi}. Timaş Yayınları.}

{\teksatir \textbf{Kitap: İki yazarlı }\par}

{\teksatir [Metin içinde atıf: (Nurmedov ve Tarhan, 2012; Svendsen ve Løber, 2020)]\par}

\kaynak{Nurmedov, S. ve Tarhan, N. (2012). \textit{Bağımlılık: Sanal veya Gerçek }(4. Baskı). Timaş Yayınları.}

\kaynak{Svendsen, S., and Løber, L. (2020). \textit{The big picture/Academic writing: The one-hour guide} (3rd digital ed.). Hans Reitzel Forlag.}

\kaynak{\textbf{Kitap: 3 veya daha fazla yazarlı  }}

\kaynak{[Metin içinde atıf: (Ülkü Arıboğan ve ark., 2008; Wolfe ve ark., 2015)]}

\kaynak{Arıboğan, D.Ü., Ortaylı, İ. ve Yavuz, H. (2008). \textit{Cumhuriyetimize Dair}. Aşiyan Yayınevi.}

\kaynak{Wolfe, J. M., Kluender, K. R., Dennis, L. M., Bartoshuk, L. M., Herz, R. S., Lederman, S. J. and Merfeld, D. M. (2015). \textit{Sensation \& perception }(4th ed.). Sinauer Associates.}

\kaynak{\textbf{Kitap Bölümü}}

\kaynak{\notrenk{FF0000}{Yazarın soyadı, Yazarın adının baş harfi. ve Yazarın soyadı, Yazarın adının baş harfi. (Yıl). Kitap bölümünün adı. }\notrenk{FF0000}{\textit{Kitabın adı}} \notrenk{FF0000}{(Baskı sayısı, Sayfa aralığı) içinde. Yayınevi.}}

\kaynak{[Metin içinde atıf: (Hovardaoğlu, 2007; Luck, 2014)]}

\kaynak{Hovardaoğlu, S. (2007). Psikolojik ölçmenin temelleri. Davranış bilimleri için araştırma teknikleri (2. Baskı, ss. 87-125) içinde. Hatipoğlu Yayınevi.}

\kaynak{Luck, S. J. (2014). A Broad Overview of the Event-Related Potential Technique. In An introduction to event related potential technique (2nd ed., pp. 1-34). The MIT Press.}

\kaynak{\textbf{Kitap: Farklı baskılar }}

\kaynak{\notrenk{FF0000}{Yazarın soyadı, Yazarın adının baş harfi. (Yıl). }\notrenk{FF0000}{\textit{Kitabın adı }}\notrenk{FF0000}{(Baskı Sayısı). Yayınevi.}}

\kaynak{[Metin içinde atıf: (Jackson, 2019; Tarhan, 2017; Wolfe ve ark., 2015)]}

\kaynak{Jackson, L. M. (2019). \textit{The psychology of prejudice: From attitudes to social action} (2nd ed.). American Psychological Association. https://doi.org/10.1037/0000168-000}

\kaynak{Tarhan, N. (2017). \textit{Toplum psikolojisi : sosyal şizofreniden toplumsal empatiye }(10. Baskı). Timaş Yayınları.}

\kaynak{Wolfe, J. M., Kluender, K. R., Dennis, L. M., Bartoshuk, L. M., Herz, R. S., Lederman, S. J. ve Merfeld, D. M. (2015). \textit{Sensation \& perception }(4th ed.). Sinauer Associates.}

\kaynak{\textbf{Kitap: Yazarı olmayan }}

\kaynak{[Metin içinde atıf: (Türkçe bilim terimleri sözlüğü, 2011)]}

\kaynak{\textit{Türkçe bilim terimleri sözlüğü : Sosyal bilimler}. (2011). Akademi Kitabevi.}

\kaynak{\textbf{Kitap: Editörlü }}

\kaynak{\notrenk{FF0000}{Editörün soyadı, editörün adının baş harfi. (Ed.). (Yıl). }\notrenk{FF0000}{\textit{Kitabın adı}} \notrenk{FF0000}{(Baskı sayısı). Yayınevi.}}

\kaynak{[Metin içinde atıf: (Hygum ve Pedersen, 2010; Kesharwani, 2020; Tanrıdağ, 2016)]}

\kaynak{Hygum, E., ve Pedersen, P. M. (Eds.). (2010). \textit{Early childhood education: Values and practices in Denmark}. Hans Reitzel Forlag.}

\kaynak{Kesharwani, P. (Ed.). (2020). \textit{Nanotechnology based approaches for tuberculosis treatment}. Academic Press}

\kaynak{Tanrıdağ, O. (Ed.). (2016). \textit{Davranış nörolojisi : Beyin-davranış ilişkilerinin organizasyon prensipleri, sendromları ve hastalıkları}. Nobel Tıp Kitapları}

\kaynak{\textbf{Editörlü Kitapta Bölüm}}

\kaynak{\notrenk{FF0000}{Yazarın soyadı, Yazarın adının baş harfi. (Yıl). Kitap bölümünün adı. Editörün adının baş harfi. Editörün soyadı (Ed.), }\notrenk{FF0000}{\textit{Kitabın adı}} \notrenk{FF0000}{(Baskı sayısı, Sayfa aralığı) içinde. Yayınevi.}}

\kaynak{[Metin içinde atıf: (Özkan Ceylan ve Bekçi, 2012; Posner, 1995)]}

\kaynak{Özkan Ceylan, A. ve Bekçi, B. (2012). Algı. N. Güngör Ergan, B. Şahin Kütük ve R. Coştur (Ed), \textit{Davranış bilimleri} (ss. 33-53) içinde. Siyasal Kitabevi}

\kaynak{Posner, M. I. (1995). Attention in cognitive neuroscience: An overview. In M. S. Gazzaniga (Ed.), \textit{The cognitive neurosciences} (pp. 615-624). The MIT Press}

\kaynak{\textbf{Kitap: Yazarı kurum olan}}

\kaynak{[Metin içinde atıf: (Üsküdar Üniversitesi, 2018)]}

\kaynak{Üsküdar Üniversitesi Yayınları. (2018). \textit{Yükseköğretim ve demokratik kültür : Vatandaşlık, insan hakları ve sivil sorumluluk, 30 Şubat 2018 }. Üsküdar Üniversitesi}

\kaynak{\textbf{Kitap: Çeviri }}

\kaynak{\notrenk{FF0000}{Orijinal kitabın yazarının soyadı, Adının baş harfi. (Yıl). }\notrenk{FF0000}{\textit{Kitabın adı}} \notrenk{FF0000}{(Baskı sayısı). (Çevirmenin adının baş harfi. Çevirmenin soyadı, Çev.). Yayınevi. (Orijinal eserin yayın tarihi )}}

\kaynak{[Metin içinde atıf: (Arasteh, 2000)]}

\kaynak{Arasteh, A. R. (2000). \textit{Aşkta ve yaratıcılıkta yeniden doğuş : Mevlana Celaleddin Rumi’nin kişilik çözümlemesi}. (Çev. B. Demirkol, ve İ. Özdemir,). Kitabiyat.}

\kaynak{\notrenk{0070C0}{\textbf{DERGİ MAKALELERİ}}}

\kaynak{\textbf{Basılı dergi makalesi }}

\kaynak{\notrenk{FF0000}{Yazarın soyadı, Yazarın Adının Baş Harfi. (Yıl). Makalenin başlığı. }\notrenk{FF0000}{\textit{Dergi adı}}\notrenk{FF0000}{\textit{ (İtalik)}}\notrenk{FF0000}{, Cilt(sayı), Sayfa aralığı. http://doi.org/xx.xxxxxxxxxx}}

\kaynak{[Metin içinde atıf: (Grady ve ark., 2019; Metin ve ark., 2018; Tarhan, 2001)}

\kaynak{Alper, E., ve Nevzat, T. (2020). Inflammation Biomarkers in Psychiatry. \textit{Current Psychiatry Research and Reviews}, \textit{16}(2), 78-85. https://doi.org/10.2174/2666082216999200625115701}

\kaynak{Grady, J. S., Her, M., Moreno, G., Perez, C., \& Yelinek, J. (2019). Emotions in storybooks: A comparison of storybooks that represent ethnic and racial groups in the United States. \textit{Psychology of Popular Media Culture}, 8(3), 207–217. https://doi.org/10.1037/ppm0000185}

\kaynak{Kalnay, E., Kanamitsu, M., Kistler, R., Collins, W., Deaven, D., Gandin, L., Iredell, M., Saha, S., White, G., Woollen, J., Zhu, Y., Chelliah, M., Ebisuzaki, W., Higgins, W., Janowiak, J., Mo, K. C., Ropelewski, C., Wang, J., Leetmaa, A., \ldots{} Joseph, D. (1996). The NCEP/NCAR 40-year reanalysis project. \textit{Bulletin of the American Meteorological Society}, 77(3), 437–471. http://doi.org/.....}

\kaynak{Metin, S. Z., Erguzel, T. T., Ertan, G., Salcini, C., Kocarslan, B., Cebi, M., Metin, B., Tanridag, O., \& Tarhan, N. (2018). The Use of Quantitative EEG for Differentiating Frontotemporal Dementia From Late-Onset Bipolar Disorder. \textit{Clinical EEG and Neuroscience}, 49(3), 171-176. https://doi.org/10.1177/1550059417750914}

\kaynak{Tarhan, N. (2001). Loneliness and Social Dissatisfaction in Turkish Adolescents. \textit{The Journal of Psychology}, \textit{135}(1), 113-123. https://doi.org/10.1080/00223980109603684}

\kaynak{\textbf{Elektronik dergi makalesi: Ön baskı }}

\kaynak{[Metin içinde atıf: (Hampton ve ark., 2017; Hetland ve ark., 20180)]}

\kaynak{Hetland, B., McAndrew, N., Perazzo, J., \& Hickman, R. (2018). A qualitative study of factors that influence active family involvement with patient care in the ICU: Survey of critical care nurse\textit{s}. \textit{PubMed Central}. https://www.ncbi.nlm.nih.gov/pmc/articles/PMC5736422/?report=classic}

\kaynak{Hampton, S., Rabagliati, H., Sorace, A., \& Fletcher-Watson, S. (2017). Autism and bilingualism: A qualitative interview study of parents’ perspectives and experience\textit{s}. \textit{PsyArXiv}. https://doi.org/10.31234/osf.io/76xfs}

\kaynak{\notrenk{0070C0}{\textbf{SÖZLÜK/ANSİKLOPEDİ}}}

\kaynak{\textbf{Sözlük ve Ansiklopedi}}

\kaynak{\notrenk{FF0000}{Yazarın soyadı, Adının baş harfi. (Yıl). }\notrenk{FF0000}{\textit{Kitabın adı}} \notrenk{FF0000}{(Baskı sayısı). Yayınevi.}}

\kaynak{[Metin içinde atıf: (Bilgin, 2003; Keyormarsi ve ark., 2020; VandenBos, 2013)]}

\kaynak{Bilgin, N. (2003). Sosyal psikoloji sözlüğü: Kavramlar, yaklaşımlar. Bağlam Yayıncılık.}

\kaynak{Keyormarsi, K., O’Leary, N., ve Pardee, A. B. (2020). \textit{Cell division}. McGraw-Hill Education. https://doi.org/10.1036/1097-8542.116300}

\kaynak{VandenBos, G. R. (Ed.). (2013). \textit{APA dictionary of clinical psychology}. American Psychological Association.}

\kaynak{\notrenk{0070C0}{\textbf{WEB KAYNAKLARI}}}

\kaynak{\textbf{Web Sayfası  }}

\kaynak{\notrenk{FF0000}{Yazar Soyadı, Adının baş harfi veya Grup adı. (Alıntı tarihi).Çalışmanın başlığı. İnternet sitesinin adı. URL adresinden .... tarihinde alınmıştır.}}

\kaynak{[Metin içinde atıf: (Türk Psikologlar Derneği, 2019; World Health Organization [WHO], 2020)]}

\kaynak{Türk Psikologlar Derneği (2019, 26 Kasım). Mesleki mevzuat. https://www.psikolog.org.tr/tr/kurumsal/meslekimevzuat-x654/ adresinden 2 Eylül 2020 tarihinde alınmıştır.}

\kaynak{World Health Organization (2020, 2 Eylül). Coronavirus. https://www.who.int/healthtopics/coronavirus\#tab=tab\_1}

\kaynak{\notrenk{0070C0}{\textbf{TEZLER}}}

\kaynak{\notrenk{FF0000}{Yazarın soyadı, Yazarın adının baş harfi. (Yıl). }\notrenk{FF0000}{\textit{Tezin baş}}\notrenk{FF0000}{lığı [Yüksek lisans tezi / Doktora tezi]. Üniversitenin adı.}}

{\teksatir [Metin içinde atıf: (Çetinkaya, 2015; Harris 2014)]\par}

\kaynak{Çetinkaya, Ş. (2015). \textit{Stochastic mortality using non-life methods} [Doktora tezi] Doğuş Üniversitesi}

\kaynak{Harris, L. (2014). \textit{Instructional leadership perceptions and practices of elementary school leaders} [Doktora tezi]. University of Virginia.}

{\teksatir Kabir, J. M. (2016). \textit{Factors influencing customer satisfaction at a fast food hamburger chain: The relationship between customer satisfaction and customer loyalty} [Doktora Tezi] Wilmington University\par}

{\teksatir Miranda, C. (2019). \textit{Exploring the lived experiences of foster youth who obtained graduate level degrees: Self-efficacy, resilience, and the impact on identity development} [Doktora tezi] Pepperdine University\par}

{\teksatir \notrenk{0070C0}{\textbf{KONFERANS BILDIRILERI }}\par}

{\teksatir \textbf{Dergide Yayınlanan Konferans Bildirileri}\par}

\kaynak{\notrenk{FF0000}{Yazarın soyadı, Yazarın Adının Baş Harfi. (Yıl). Bildirinin başlığı. }\notrenk{FF0000}{\textit{Derginin Adı}}\notrenk{FF0000}{, Cilt(Süreli yayının sayısı), Sayfa aralığı. http://doi.org/xx.xxxxxxxxxx}}

\kaynak{[Metin içinde atıf: (Duckworth ve ark., 2019)]}

\kaynak{Duckworth, A. L., Quirk, A., Gallop, R., Hoyle, R. H., Kelly, D. R., ve Matthews, M. D. (2019). Cognitive and noncognitive predictors of success. \textit{Proceedings of the National Academy of Sciences}, USA, 116(47), 23499–23504. https://doi.org/10.1073/pnas.1910510116}

{\teksatir \textbf{Yayımlanmamış konferans bildirisi}\par}

\kaynak{\notrenk{FF0000}{Yazarın soyadı, Yazarın adının baş harfi. (Yıl, Gün, Ay). Sözlü/poster sunumun adı [Sözlü/poster Sunum]. Bilimsel Toplantının Adı, Toplantının gerçekleştiği şehir, Ülke.}}

\kaynak{Özkan A. ve Şenyüz L. (2004, Eylül, 7-11). Haloperidolün, farklı uyaranlar tarafından kontrol edilen, su elde etmeye yönelik edimsel davranış üzerindeki etkileri [Sözlü sunum]. 13. Ulusal Psikoloji Kongresi, İstanbul, Türkiye.}

\kaynak{Salman, F. ve Cangöz, B. (2016, Eylül, 5-7). Alzheimer tipi demansta iç kaynaklı ve dış kaynaklı dikkat: duygusal bağlam etkisi [Poster sunumu]. 19. Ulusal Psikoloji Kongresi, İzmir, Türkiye.}

{\teksatir \textbf{Basılmış/Elektronik konferans kitabında bildiri }\par}

\kaynak{\notrenk{FF0000}{Yazarın soyadı, Yazarın adının baş harfi. (Yıl, Gün, Ay). Bildirinin adı. Editörün adının baş harfi. Editörün soyadı (Ed.), }\notrenk{FF0000}{\textit{Sempozyum Kitabının adı}} \notrenk{FF0000}{(Baskı sayısı, Sayfa aralığı) içinde. Toplantının gerçekleştiği şehir, Ülke}}

\kaynak{Altürk, S., Avci, D., Tamer, Ö. ve Atalay, Y. (2017, 6-10, Eylül). The second− and third−order nonlinear optical properties and Electronic transitions of a nlo chromophore: a DFTstudy. B. Akkuş, Y. Öktem, G. Süsoy, Ö. Aytan, A. Kurt, G. Aydoğdu, O. A.Barut (Ed.), \textit{9 th International Physics Conference of the Balkan Physical Union} içinde. İstanbul, Türkiye}

{\teksatir \notrenk{0070C0}{\textbf{RAPORLAR }}\par}

{\teksatir \textbf{Resmi raporlar }\par}

{\teksatir \notrenk{FF0000}{Yazarın soyadı, Yazarın Adının Baş Harfi. (Yıl). }\notrenk{FF0000}{\textit{Raporun}}\notrenk{FF0000}{\textit{ }}\notrenk{FF0000}{\textit{adı}} \notrenk{FF0000}{(Yayın no.).Yayınevi/Erişim adresi}\par}

{\teksatir \notrenk{FF0000}{Resmi Yayını Basan Kurumun Adı. (Yıl). }\notrenk{FF0000}{\textit{Raporun}}\notrenk{FF0000}{\textit{ }}\notrenk{FF0000}{\textit{adı}} \notrenk{FF0000}{(Yayın no.).Yayınevi/Erişim adresi}\par}

{\teksatir [Metin içinde atıf: (Akbaytürk Çanak ve ark., 2014, s. 40; Türkiye İstatistik Kurumu [TÜİK], 2019, s. 67)]\par}

\kaynak{Akbaytürk Çanak, T., Çelik, S., Çetinkaya, İ., Çukadar, S., Güneş, G., Gürdal, G., Holt İ., Kaya E. ve Kaygusuz, A. (2014). \textit{2023’e doğru Türkiye’de üniversite kütüphaneleri: Mevcut durum, sorunlar, standartlar ve çözüm önerileri}. Yükseköğretim Kurulu. Erişim adresi https://yok.gov.tr/web/ukr}

\kaynak{Türkiye İstatistik Kurumu. (2019). \textit{İstatistiklerle çocuk} (Yayın no.4581). https://biruni.tuik.gov.tr/yayin/views/visitorPages/index.zul}

{\teksatir \notrenk{0070C0}{\textbf{PATENT }}\par}

{\teksatir \notrenk{FF0000}{Yazarın soyadı, Yazarın Adının Baş Harfi. (Yıl). }\notrenk{FF0000}{\textit{Patentin}}\notrenk{FF0000}{\textit{ }}\notrenk{FF0000}{\textit{adı}}\notrenk{FF0000}{. Patent Ülkesi. Patent No}\par}

{\teksatir [Metin içinde atıf: (Apelsmeier, 2013)]\par}

{\teksatir Apelsmeier, A. (2013).\textit{ Motor vehicle having an air conditioner compressor as a starter of the internal combustion engine. }U.S. Patent No. US9731704B2\textit{ }\par}

\begin{kaynaklistesi}
  \item Abrahart, R. J. ve See, L. (1998, 17–19, Eylül). Neural Network vs. ARMA Modelling: Constructing Benchmark Case Studies of River Flow Prediction. In GeoComputation ’98. Proceedings of the Third International Conference on GeoComputation, University of Bristol, United Kingdom

  \item Abrahart, R. J. ve See, L. (2000). Comparing neural network and autoregressive moving average techniques for the provision of continuous river flow forecasts in two contrasting catchments, \textit{Hydrological Processes,} 14, 2157–2172.

  \item Spielberg, S. (Producer) ve Spielberg, S. (Director). (1993). \textit{Schindler’s list }[Motion picture]. California : Universal Picture.

\end{kaynaklistesi}

